\chapter{Instrukcja uruchomienia}
\label{cha:chapter005}

\section{Kroki uruchomienia programu}

Poniższy przewodnik umożliwia samodzielne uruchomienie programu od
konfiguracji środowiska po uruchomienie inferencji na żywo.

\subsection*{Krok 1: Pobranie repozytorium}

\begin{lstlisting}[style=styleBash, numbers=none, caption={Klonowanie repozytorium}, label=lst:clone]
git clone https://github.com/JustFiesta/mnist-translator.git
cd mnist-translator
\end{lstlisting}

\subsection*{Krok 2: Instalacja środowiska Python}

\begin{lstlisting}[style=styleBash, numbers=none, caption={Konfiguracja środowiska wirtualnego}, label=lst:venv]
# Instalacja Pythona 3.13 przez menedzera uv
uv python install 3.13
uv venv --python 3.13

# Windows PowerShell
.venv\Scripts\Activate.ps1

# Linux / macOS
source .venv/bin/activate

# Instalacja zaleznosci z pliku lock
uv sync
\end{lstlisting}

\subsection*{Krok 3: Pobranie zbioru danych}

\begin{enumerate}
    \item Przejdź na stronę Kaggle:\\
          \url{https://www.kaggle.com/datasets/datamunge/sign-language-mnist/data}
    \item Pobierz pliki \texttt{sign\_mnist\_train.csv} i \texttt{sign\_mnist\_test.csv}.
    \item Utwórz katalog docelowy i umieść w nim pobrane pliki:
\end{enumerate}

\begin{lstlisting}[style=styleBash, numbers=none, caption={Przygotowanie struktury katalogów na dane}, label=lst:dirs]
# Windows PowerShell
New-Item -ItemType Directory -Path data/raw -Force

# Linux / macOS
mkdir -p data/raw

# Skopiuj pobrane pliki CSV do katalogu data/raw/
\end{lstlisting}

\subsection*{Krok 4: Trenowanie modeli}

\begin{lstlisting}[style=styleBash, numbers=none, caption={Trenowanie wszystkich trzech klasyfikatorów}, label=lst:train-all]
# SVC (domyslny, najszybszy)
uv run python main.py train

# Las Losowy
uv run python main.py train --classifier rf --model artifacts/rf_model.pkl

# CNN (trenowanie moze zajac kilka minut)
uv run python main.py train --classifier cnn --model artifacts/cnn_model.keras
\end{lstlisting}

\subsection*{Krok 5: Ewaluacja modeli}

\begin{lstlisting}[style=styleBash, numbers=none, caption={Ewaluacja wytrenowanych modeli}, label=lst:eval-all]
# Ewaluacja SVC
uv run python main.py eval \
    --model artifacts/svc_model.pkl \
    --dataset data/raw/sign_mnist_test.csv

# Ewaluacja Lasu Losowego
uv run python main.py eval \
    --model artifacts/rf_model.pkl \
    --dataset data/raw/sign_mnist_test.csv

# Ewaluacja CNN
uv run python main.py eval \
    --model artifacts/cnn_model.keras \
    --dataset data/raw/sign_mnist_test.csv
\end{lstlisting}

Zanotuj dokładność każdego modelu i porównaj z wynikami z tabeli~\ref{tab:wyniki}.

\subsection*{Krok 6: Inferencja offline (z pliku CSV)}

\begin{lstlisting}[style=styleBash, numbers=none, caption={Inferencja na wybranym wierszu pliku CSV}, label=lst:infer-csv]
uv run python main.py infer \
    --source csv \
    --row 42 \
    --dataset data/raw/sign_mnist_test.csv \
    --model artifacts/svc_model.pkl
\end{lstlisting}

\newpage
\subsection*{Krok 7: Inferencja na żywo z kamery}

\begin{lstlisting}[style=styleBash, numbers=none, caption={Uruchomienie inferencji w czasie rzeczywistym}, label=lst:infer-live]
uv run python main.py infer \
    --source camera \
    --camera 0 \
    --model artifacts/svc_model.pkl
\end{lstlisting}

\textbf{Uwagi:}
\begin{itemize}
    \item Przy pierwszym uruchomieniu aplikacja automatycznie pobierze plik
          modelu MediaPipe (\texttt{gesture\_recognizer.task}) --- wymagany
          jest dostęp do Internetu.
    \item Naciśnij klawisz \texttt{q}, aby zakończyć działanie okna kamery.
    \item Jeśli kamera nie jest wykrywana pod indeksem 0, spróbuj
          opcji \texttt{-{}-camera 1}.
\end{itemize}

\subsection*{Krok 8: Uruchomienie testów jednostkowych}

\begin{lstlisting}[style=styleBash, numbers=none, caption={Uruchamianie testów pytest}, label=lst:tests]
# Wszystkie testy
uv run pytest

# Testy z raportem pokrycia kodu
uv run pytest --cov=src --cov-report=term-missing
\end{lstlisting}

